% Einheit 1 – Terme aufstellen und berechnen (bank/terme/e1.jsonl, zusammenbau v1.0)
\einheitenkopf[e1][Terme aufstellen und berechnen]{Terme aufstellen und berechnen}
\setcounter{aufgabe}{26}


% test: terme-e1-k1-s7-v2
\begin{kbaufgabe}{Kannst du das schon? Dann bist du fertig}
\begin{kbblock}
\kbteil{}{Die Summe aus dem Vierfachen einer Zahl $x$ und 9 wird halbiert. Kreuze den passenden Term an.}{P10 ’23}
\kbkreuz{$4x + 9 : 2$}
\kbkreuz{$(4x + 9) : 2$}
\kbkreuz{$2 \cdot (4x + 9)$}
\kbkreuz{$4x : 9 + 2$}
\end{kbblock}
\end{kbaufgabe}


% ohne: terme-e1-k2-s1-v1, terme-e1-k2-s1-v2, terme-e1-k2-s1-v3
\begin{kbaufgabe}{Termwerte berechnen}
\begin{kbblock}
\kbteil{a)}{Berechne den Wert des Terms $4x - 3$ für $x = 5$.}{}
\kbantwort{\kbfeld}
\end{kbblock}
\begin{kbblock}
\kbteil{b)}{Berechne den Wert des Terms $2x + 9$ für $x = -3$.}{}
\kbantwort{\kbfeld}
\end{kbblock}
\begin{kbblock}
\kbteil{c)}{Berechne den Wert des Terms $10 - 3x$ für $x = -2$.}{}
\kbantwort{\kbfeld}
\end{kbblock}
\end{kbaufgabe}


% leiter: terme-e1-k1-s1-v1, terme-e1-k1-s1-v2, terme-e1-k1-s1-v3, terme-e1-k1-s2-v1, terme-e1-k1-s3-v1, terme-e1-k1-s5-v1, terme-e1-k1-s6-v1
\begin{kbaufgabe}{Einen Term aufstellen}
\begin{kbblock}
\kbteil{a)}{Das Vierfache einer Zahl $x$ wird um 3 vermindert. Kreuze den Term an, der dazu passt.}{}
\kbkreuz{$4x - 3$}
\kbkreuz{$4 \cdot (x - 3)$}
\kbkreuz{$3 - 4x$}
\end{kbblock}
\begin{kbblock}
\kbteil{b)}{Das Vierfache einer Zahl $x$ wird um 5 vermindert. Kreuze den Term an, der dazu passt.}{}
\kbkreuz{$5 - 4x$}
\kbkreuz{$4x - 5$}
\kbkreuz{$4 \cdot (x - 5)$}
\end{kbblock}
\begin{kbblock}
\kbteil{c)}{Das Vierfache einer Zahl $x$ wird um 6 vermindert. Kreuze den Term an, der dazu passt.}{}
\kbkreuz{$4 \cdot (x - 6)$}
\kbkreuz{$6 - 4x$}
\kbkreuz{$4x - 6$}
\end{kbblock}
\begin{kbblock}
\kbteil{d)}{Das Dreifache einer Zahl $x$ wird um 8 vergrößert. Schreibe dazu einen Term.}{}
\kbantwort{\kbfeld}
\end{kbblock}
\begin{kbblock}
\kbteil{e)}{Die Figur zeigt ein Rechteck. Die Maße sind in cm angegeben. Schreibe einen Term für den Flächeninhalt des Rechtecks.}{}
\kbzweispaltig[0.42]{\viereck[seiten={5,x,5,x}]{(0,0)}{(5,0)}{(5,2.5)}{(0,2.5)}}{\lbantwortrechts[1.5cm]{\kbfeld}}
\end{kbblock}
\begin{kbblock}
\kbteil{f)}{Ein Rechteck hat die Seiten $x$ cm und 9 cm. Schreibe einen Term für den Umfang. Fasse den Term zusammen.}{}
\kbantwort{\kbfeld}
\end{kbblock}
\begin{kbblock}
\kbteil{g)}{Eine Rechnung geht so: Verdreifache eine Zahl $x$, addiere 2 und verdopple die Summe. Schreibe dazu einen Term.}{}
\kbantwort{\kbfeld}
\end{kbblock}
\end{kbaufgabe}


% pruefung: terme-e1-k1-s7-v1, terme-e1-k1-s7-v3
\begin{kbaufgabe}{Den passenden Term auch in Aufgaben aus der Prüfung finden}
\begin{kbblock}
\kbteil{a)}{Die Differenz aus dem Dreifachen einer Zahl $x$ und 7 wird verdoppelt. Kreuze den passenden Term an.}{P10 ’23}
\kbkreuz{$(3x - 7) : 2$}
\kbkreuz{$3x : 7 \cdot 2$}
\kbkreuz{$2 \cdot (3x - 7)$}
\kbkreuz{$3x - 7 \cdot 2$}
\end{kbblock}
\begin{kbblock}
\kbteil{b)}{Das Sechsfache einer Zahl $x$ wird um zwei vermindert. Kreuze den passenden Term an.}{P10 ’17}
\kbkreuz{$6 \cdot (x - 2)$}
\kbkreuz{$2 - 6x$}
\kbkreuz{$6x - 2$}
\kbkreuz{$6x - 2x$}
\end{kbblock}
\end{kbaufgabe}


% ohne: terme-e1-k3-s1-v1, terme-e1-k3-s1-v2, terme-e1-k3-s1-v3
\begin{kbaufgabe}{Zu einem Term eine passende Situation erfinden}
\begin{kbblock}
\kbteil{a)}{Erfinde eine Situation aus dem Alltag, zu der der Term $15 + 3x$ passt. Schreibe auf, wofür $x$ steht. Berechne den Term für $x = 2$ und erkläre, was das Ergebnis bedeutet.}{}
\item[]\kbraum{2}
\end{kbblock}
\begin{kbblock}
\kbteil{b)}{Erfinde eine Situation aus dem Alltag, zu der der Term $20 - 2x$ passt. Schreibe auf, wofür $x$ steht. Berechne den Term für $x = 4$ und erkläre, was das Ergebnis bedeutet.}{}
\item[]\kbraum{2}
\end{kbblock}
\begin{kbblock}
\kbteil{c)}{Erfinde eine Situation aus dem Alltag, zu der der Term $10 + 4x$ passt. Schreibe auf, wofür $x$ steht. Berechne den Term für $x = 3$ und erkläre, was das Ergebnis bedeutet.}{}
\item[]\kbraum{2}
\end{kbblock}
\end{kbaufgabe}


% pflicht: terme-e1-k4-s1-v1, terme-e1-k4-s2-v1, terme-e1-k4-s3-v1, terme-e1-k4-s4-v1
\begin{kbaufgabe}{Verstanden?}
\begin{kbblock}
\kbteil{a)}{Mia soll zu „Eine Zahl $x$ wird um 9 vermindert“ einen Term schreiben. Sie schreibt: $9 - x$. Finde den Fehler und schreibe den Term richtig.}{}
\item[]\kbraum{2}
\end{kbblock}
\begin{kbblock}
\kbteil{b)}{Begründe, warum man für „das Doppelte der Summe aus $x$ und 5“ eine Klammer braucht. Setze dazu $x = 3$ in $2 \cdot (x + 5)$ und in $2x + 5$ ein.}{}
\item[]\kbraum{3}
\end{kbblock}
\begin{kbblock}
\kbteil{c)}{Welcher Term gehört zum Termbaum?}{}
\kbzweispaltig[0.38]{\termbaum{\cdot}{4}{\tb{+}{x}{6}}}{\lbantwortrechts[1.1cm]{\kbfeld}}
\end{kbblock}
\begin{kbblock}
\kbteil{d)}{Ein Stromtarif kostet 11 € Grundpreis im Monat und 0,30 € je Kilowattstunde. Stelle einen Term für die Monatskosten in Euro bei $n$ Kilowattstunden auf. Was kostet ein Monat mit 150 Kilowattstunden?}{}
\item[]\kbraum{3}
\end{kbblock}
\end{kbaufgabe}
